% ═══════════════════════════════════════════════════════════════
% LEHRERHINWEISE: Reaktionsgleichungen aufstellen
% Thema 2.4b | Chemie Klasse 8 | inkl. Übung exotische Oxide
% ═══════════════════════════════════════════════════════════════

\documentclass[11pt, a4paper]{article}

\usepackage[utf8]{inputenc}
\usepackage[T1]{fontenc}
\usepackage[ngerman]{babel}

\usepackage{helvet}
\renewcommand{\familydefault}{\sfdefault}

\usepackage{microtype}
\usepackage[margin=2cm, top=2cm, bottom=2cm]{geometry}
\usepackage{enumitem}
\usepackage{tabularx}
\usepackage{booktabs}
\usepackage{array}
\usepackage{multicol}

\usepackage{fancyhdr}
\usepackage{lastpage}

\usepackage{xcolor}
\usepackage{tcolorbox}
\tcbuselibrary{skins}

% ═══════════════════════════════════════════════════════════════
% FARBEN
% ═══════════════════════════════════════════════════════════════

\definecolor{chemie}{HTML}{2a7a4b}
\definecolor{hellgrau}{HTML}{F5F5F5}
\definecolor{mittelgrau}{HTML}{888888}
\definecolor{dunkelgrau}{HTML}{444444}
\definecolor{loesung}{HTML}{C62828}

\colorlet{fachfarbe}{chemie}

% ═══════════════════════════════════════════════════════════════
% KOPF- UND FUSSZEILE
% ═══════════════════════════════════════════════════════════════

\pagestyle{fancy}
\fancyhf{}
\renewcommand{\headrulewidth}{0.4pt}
\renewcommand{\footrulewidth}{0pt}
\lhead{\small Chemie Klasse 8 | Thema 2.4b}
\rhead{\small\textcolor{fachfarbe}{\textbf{LEHRERHINWEISE}}}
\cfoot{\small\textcolor{mittelgrau}{Seite \thepage{} von \pageref{LastPage}}}

% ═══════════════════════════════════════════════════════════════
% BOXEN
% ═══════════════════════════════════════════════════════════════

\newtcolorbox{infobox}[1][]{
    colback=hellgrau,
    colframe=hellgrau,
    boxrule=0pt,
    arc=0pt,
    left=8pt, right=8pt, top=6pt, bottom=6pt,
    borderline west={3pt}{0pt}{fachfarbe},
    #1
}

\newtcolorbox{warnbox}{
    colback=white,
    colframe=loesung,
    boxrule=1pt,
    arc=0pt,
    left=8pt, right=8pt, top=6pt, bottom=6pt
}

% ═══════════════════════════════════════════════════════════════
% BEFEHLE
% ═══════════════════════════════════════════════════════════════

\newcommand{\abschnitt}[1]{%
    \vspace{10pt}%
    \noindent{\large\bfseries\textcolor{fachfarbe}{#1}}%
    \vspace{6pt}%
    \nopagebreak[4]%
}

\newcommand{\lsg}[1]{\textcolor{loesung}{\textbf{#1}}}

\setlength{\parindent}{0pt}
\setlength{\parskip}{4pt}

\newcolumntype{L}[1]{>{\raggedright\arraybackslash}p{#1}}
\newcolumntype{C}[1]{>{\centering\arraybackslash}p{#1}}

% ═══════════════════════════════════════════════════════════════
% DOKUMENT
% ═══════════════════════════════════════════════════════════════

\begin{document}

% ═══════════════════════════════════════════════════════════════
% HEADER
% ═══════════════════════════════════════════════════════════════

\begin{center}
    {\Large\bfseries\textcolor{fachfarbe}{Reaktionsgleichungen aufstellen}}\\[0.3em]
    {\normalsize Lehrerhinweise | Chemie Klasse 8 | Thema 2.4b}
\end{center}
\vspace{-0.3em}
\hrule
\vspace{0.8em}

% ═══════════════════════════════════════════════════════════════
% ÜBERBLICK
% ═══════════════════════════════════════════════════════════════

\abschnitt{Überblick}

\begin{infobox}
\textbf{Lernziele:}
\begin{itemize}[leftmargin=1.5em, itemsep=0pt]
    \item SuS wenden die Kreuzregel zur Formelermittlung an
    \item SuS stellen Reaktionsgleichungen für Oxidationen auf
    \item SuS gleichen Reaktionsgleichungen korrekt aus
    \item SuS unterscheiden Index und Vorfaktor
\end{itemize}
\end{infobox}

\textbf{Voraussetzungen (aus DS 1):}
\begin{itemize}[leftmargin=1.5em, itemsep=0pt]
    \item Oxidation als Reaktion mit Sauerstoff
    \item Wortgleichungen formulieren
    \item Wertigkeiten der Hauptgruppenelemente
\end{itemize}

\textbf{Materialien:}
\begin{itemize}[leftmargin=1.5em, itemsep=0pt]
    \item AB-02-reaktionsgleichungen (Grundlagen)
    \item UEBUNG-02b-exotische-oxide (Vertiefung, 1× pro SuS oder Tafel)
    \item ML-02b-exotische-oxide (Musterlösung mit Schritt-für-Schritt)
    \item Tablets für LP-02 und MT-02
\end{itemize}

\textbf{Digitale Ressourcen:}
\begin{itemize}[leftmargin=1.5em, itemsep=0pt]
    \item \texttt{LP-02-reaktionsgleichungen.html} -- Interaktiver Lernpfad
    \item \texttt{MT-02-reaktionsgleichungen.html} -- Stundenleistung (33/48 BE)
\end{itemize}

% ═══════════════════════════════════════════════════════════════
% ÜBUNG EXOTISCHE OXIDE
% ═══════════════════════════════════════════════════════════════

\abschnitt{Übung: Exotische Oxide (UEBUNG-02b)}

\textbf{Ziel:} Festigung der Kreuzregel und Reaktionsgleichungen vor Einführung der Reduktion

\textbf{Didaktischer Hinweis:} Silberoxid (Ag$_2$O) steht bewusst am Ende als Überleitung zur Reduktion -- die Schüler werden später sehen, wie diese Reaktion umgekehrt wird.

\begin{center}
\small
\renewcommand{\arraystretch}{1.4}
\begin{tabular}{|C{0.5cm}|L{2.8cm}|C{1.5cm}|C{2cm}|L{5.5cm}|}
\hline
\textbf{\#} & \textbf{Oxid} & \textbf{Wert.} & \textbf{Formel} & \textbf{Reaktionsgleichung} \\
\hline
1 & Bariumoxid & II & \lsg{BaO} & \lsg{2 Ba + O$_2$ $\rightarrow$ 2 BaO} \\
\hline
2 & Natriumoxid & I & \lsg{Na$_2$O} & \lsg{4 Na + O$_2$ $\rightarrow$ 2 Na$_2$O} \\
\hline
3 & Chromoxid & III & \lsg{Cr$_2$O$_3$} & \lsg{4 Cr + 3 O$_2$ $\rightarrow$ 2 Cr$_2$O$_3$} \\
\hline
4 & Titandioxid & IV & \lsg{TiO$_2$} & \lsg{Ti + O$_2$ $\rightarrow$ TiO$_2$} \\
\hline
5 & Manganoxid & IV & \lsg{MnO$_2$} & \lsg{Mn + O$_2$ $\rightarrow$ MnO$_2$} \\
\hline
6 & Zinnoxid & IV & \lsg{SnO$_2$} & \lsg{Sn + O$_2$ $\rightarrow$ SnO$_2$} \\
\hline
7 & Silberoxid & I & \lsg{Ag$_2$O} & \lsg{4 Ag + O$_2$ $\rightarrow$ 2 Ag$_2$O} \\
\hline
\end{tabular}
\end{center}

\textbf{Überleitung zur Reduktion:}\\
``Silberoxid ist das, was angelaufenes Silber schwarz macht. Kann man diese Reaktion umkehren? Kann man aus dem schwarzen Silberoxid wieder glänzendes Silber gewinnen?''\\
$\rightarrow$ Nächstes Thema: \textbf{Reduktion} (Zerlegung von Ag$_2$O durch Erhitzen)

% ═══════════════════════════════════════════════════════════════
% STUNDENVERLAUF
% ═══════════════════════════════════════════════════════════════

\abschnitt{Stundenverlauf: Wiederholung + Reduktion}

\begin{center}
\small
\begin{tabular}{|C{1.2cm}|L{2.5cm}|L{6cm}|L{3.5cm}|}
\hline
\textbf{Min} & \textbf{Phase} & \textbf{Inhalt} & \textbf{Material} \\
\hline
0--5 & Einstieg & Kurze Wiederholung: Was ist Oxidation? Kreuzregel? & Tafel \\
\hline
5--15 & Übung I & UEBUNG-02b Aufgabe A: Oxidformeln (\#1--3 gemeinsam, \#4--7 EA) & UEBUNG-02b \\
\hline
15--25 & Übung II & UEBUNG-02b Aufgabe B: Reaktionsgleichungen (Tafel: \#2, \#3) & Tafel, ML-02b \\
\hline
25--30 & Überleitung & \#7 Silberoxid gemeinsam; Frage: ``Kann man das umkehren?'' & Tafel \\
\hline
\multicolumn{4}{|c|}{\textit{--- Pause / Themenwechsel ---}} \\
\hline
30--90 & Reduktion & LP-03, AB-03, SIM-03 (siehe PLANUNG-03-reduktion.md) & LP-03, AB-03 \\
\hline
\end{tabular}
\end{center}

% ═══════════════════════════════════════════════════════════════
% HÄUFIGE FEHLER
% ═══════════════════════════════════════════════════════════════

\abschnitt{Häufige Fehler und Reaktionen}

\begin{warnbox}
\textbf{1. Kreuzregel ohne Kürzen}\\
\textit{Fehler:} Ba (II) + O (II) $\rightarrow$ Ba$_2$O$_2$\\
\textit{Reaktion:} ``Wenn beide Wertigkeiten gleich sind, kürzen wir auf 1:1''

\vspace{6pt}

\textbf{2. Sauerstoff als Atom statt Molekül}\\
\textit{Fehler:} 2 Ba + 2 O $\rightarrow$ 2 BaO\\
\textit{Reaktion:} ``Sauerstoff kommt in der Natur \textbf{immer} als O$_2$ vor!''

\vspace{6pt}

\textbf{3. Ungerade Sauerstoffzahl im Produkt}\\
\textit{Problem:} Cr$_2$O$_3$ hat 3 O, aber O$_2$ liefert immer gerade Anzahl\\
\textit{Lösung:} Produkt verdoppeln: 2 Cr$_2$O$_3$ (6 O) $\rightarrow$ 3 O$_2$

\vspace{6pt}

\textbf{4. Index statt Vorfaktor ändern}\\
\textit{Fehler:} Ba + O$_2$ $\rightarrow$ BaO$_2$ (Index geändert!)\\
\textit{Reaktion:} ``Der Index gehört zur Formel und darf \textbf{nie} geändert werden. Nur die Vorfaktoren vor der Formel dürfen angepasst werden.''
\end{warnbox}

% ═══════════════════════════════════════════════════════════════
% DIFFERENZIERUNG
% ═══════════════════════════════════════════════════════════════

\abschnitt{Differenzierung}

\begin{tabular}{|c|L{4cm}|L{8.5cm}|}
\hline
\textbf{Niveau} & \textbf{Zielgruppe} & \textbf{Anpassung} \\
\hline
$\star$ & Berufsreife (BR) & Nur \#1, \#4--6 (einfache Wertigkeiten); Kreuzregel vorgeben; Ausgleichen mit Hilfe \\
\hline
$\star\star$ & Mittlere Reife (MR) & Alle 7 Oxide; Unterstützung bei \#3 (Chromoxid) \\
\hline
$\star\star\star$ & Gymnasium (GY) & Zusatz: Weitere Oxide selbst recherchieren (z.B. V$_2$O$_5$, WO$_3$) \\
\hline
\end{tabular}

% ═══════════════════════════════════════════════════════════════
% MT-02
% ═══════════════════════════════════════════════════════════════

\abschnitt{Stundenleistung MT-02}

\begin{infobox}
\textbf{URL:} \texttt{https://devbox42.github.io/sciencesim/chemie/kl08/02-reaktionsgleichungen/MT-02-reaktionsgleichungen.html}

\vspace{0.3em}
\textbf{Bewertung:}
\begin{tabular}{lll}
$\star$ Basis & 33 BE & Kreuzregel, einfache Gleichungen \\
$\star\star\star$ Erweitert & 48 BE & + anspruchsvolle Ausgleichsaufgaben
\end{tabular}

\vspace{0.3em}
\textbf{Hilfsmittel:} AB-02 erlaubt (Wertigkeitstabelle)\\
\textbf{Zeit:} ca. 25--30 Minuten\\
\textbf{Abgabe:} PDF speichern $\rightarrow$ Bildungscloud
\end{infobox}

% ═══════════════════════════════════════════════════════════════
% HINTERGRUND
% ═══════════════════════════════════════════════════════════════

\abschnitt{Hintergrundwissen: Die exotischen Oxide}

\small
\begin{tabular}{L{2.5cm}L{11.5cm}}
\textbf{Bariumoxid} & Feuerwerk (grüne Flammenfarbe bei Ba-Verbindungen), Trocknungsmittel, früher in Bildröhren \\[0.3em]
\textbf{Natriumoxid} & Glasherstellung (Soda-Kalk-Glas), reagiert heftig mit Wasser zu NaOH \\[0.3em]
\textbf{Chromoxid} & Pigment ``Chromgrün'' (Farben, Keramik), Schleifmittel, Katalysator; toxisch! \\[0.3em]
\textbf{Titandioxid} & Weißpigment Nr. 1 weltweit: Wandfarbe, Kunststoff, Papier, Sonnencreme (UV-Schutz), Zahnpasta, Lebensmittelfarbe E171 (umstritten) \\[0.3em]
\textbf{Manganoxid} & ``Braunstein'' in Batterien (Zink-Kohle, Alkaline), Oxidationsmittel, Glasfärbung \\[0.3em]
\textbf{Zinnoxid} & Keramik-Glasuren (weiß, deckend), Glas-Polierung, Gas-Sensoren \\[0.3em]
\textbf{Silberoxid} & Angelaufenes Silber (schwarze Schicht), Knopfzellen, lässt sich durch Erhitzen leicht zerlegen $\rightarrow$ \textbf{Reduktion} \\
\end{tabular}

\end{document}
